\begin{DefinitionBox}{Geradengleichung}{real_functions_affine}
  Eine nicht senkrechte Gerade ist der Graph einer Funktion
  $g:\R\to\R$. Ihre \emph{Geradengleichung}
  hat mit $m,b\in\R$ die Form
  \par\smallskip
  \begin{minipage}[c]{0.47\linewidth}
    {\centering
      \CachedTikzFigure{content/resources/chapters/real_functions/tikz/functions_affine_formula_annotations.tex}\par}
    \smallskip\noindent\raggedright
    \begin{itemize}[leftmargin=1.3em,itemsep=0.4em,topsep=0.3em]
      \item Am \textcolor{THMBlue}{\emph{Steigungsdreieck}}
        \EnglishTerm{slope triangle} liest man die \emph{Steigung}
        \EnglishTerm{slope} ab.
        \[
          \textcolor{THMBlue}{m=\Delta y/\Delta x}
        \]
      \item Der \textcolor{green!55!black}{$y$-Achsenabschnitt ist $b=g(0)$}.
    \end{itemize}
  \end{minipage}\hfill
  \begin{minipage}[c]{0.51\linewidth}
    \centering
    \CachedTikzFigure{content/resources/chapters/real_functions/tikz/functions_affine_line_definition.tex}
  \end{minipage}
\end{DefinitionBox}
Eine Geradengleichung dieser Form heißt auch \emph{affine Funktion} \EnglishTerm{affine function}.


Das Vorzeichen der Steigung $m$ bestimmt die Richtung der Geraden:
\[
  \begin{array}{c@{\qquad}c@{\qquad}c}
    m>0&m=0&m<0\\
    \text{steigende Gerade}&\text{waagerechte Gerade}&\text{fallende Gerade}
  \end{array}
\]

\subsection{Nullstelle einer Geraden}
\label{subsec:real_functions_line_zero}
\begin{KeyPointBox}{Nullstelle einer Geraden}{real_functions_zeros}
  Für eine Gerade $g(x)=mx+b$ mit $m\ne0$ berechnen wir die Nullstelle aus
  \[
    g(x_0)\overset{!}{=}0
    \qquad \Longleftrightarrow \qquad
    mx_0+b=0
    \qquad \Longleftrightarrow \qquad
    x_0=-\frac{b}{m}.
  \]
  Für $m=0$ liegt die Gerade entweder auf der $x$-Achse ($b=0$) oder parallel ($b\ne0$) dazu.
\end{KeyPointBox}
\begin{ExampleBox}{Nullstelle einer Geraden}{real_functions_zeros_example}
  \begin{minipage}[c]{0.50\linewidth}
    Wir betrachten die Geradengleichung 
    \[
      g(x)=\frac12x+1 \quad \text{d.h.}\quad m=\Cloze{1cm}{\frac12},\quad b=\Cloze{1cm}{1}.
    \]
    Für die Nullstelle $x_0$ ergibt sich damit
    \[
      x_0=\Cloze{6cm}{-\frac{b}{m}=-\frac{1}{\frac12}=-2}
    \]
  \end{minipage}\hfill
  \begin{minipage}[c]{0.49\linewidth}
    \raggedleft
    \CachedTikzFigure{content/resources/chapters/real_functions/tikz/functions_line_intercepts.tex}
  \end{minipage}
  Die berechnete Nullstelle bestätigt sich am Graphen.
\end{ExampleBox}

\subsection{Eine Gerade aus zwei Punkten bestimmen}
\label{subsec:real_functions_two_points}
Wie bestimmen wir die Geradengleichung durch zwei gegebene Punkte? Wir setzen
ihre Koordinaten in $g(x)=mx+b$ ein und lösen das entstehende lineare
Gleichungssystem nach $m$ und $b$.
\[
  \begin{aligned}
    mx_A+b&=y_A,\\
    mx_B+b&=y_B
  \end{aligned}
  \quad\Longleftrightarrow\quad
  \begin{aligned}
    1 b+x_A m&=y_A,\\
    1 b+x_B m&=y_B
  \end{aligned}
  \quad\Longleftrightarrow\quad
  \left[\begin{array}{@{}cc|c@{}}
    \Cloze{0.7cm}{1}&\Cloze{0.8cm}{x_A}&\Cloze{0.8cm}{y_A}\\[0.3em]
    \Cloze{0.7cm}{1}&\Cloze{0.8cm}{x_B}&\Cloze{0.8cm}{y_B}
  \end{array}\right].
\]
Wir lösen das LGS mit dem Gauß-Algorithmus. 
\par\noindent\textbf{Vorwärtselimination:} Die Operation steht rechts neben der Matrix,
auf die sie angewendet wird:
\begin{ClozeArea}[34]{7}
  \[
    \begin{array}{r@{\;}r@{\qquad}l}
      &\left[\begin{array}{@{}cc|c@{}}
        1&x_A&y_A\\
        1&x_B&y_B
      \end{array}\right]
      &\raisebox{-0.65em}{\(\TransformOp{Z_2-Z_1}\)}\\[1.2em]
      \Longleftrightarrow
      &\left[\begin{array}{@{}cc|c@{}}
        1&x_A&y_A\\
        \textcolor{THMBlue}{0}&x_B-x_A&y_B-y_A
      \end{array}\right]
      &\text{Zeilenstufenform}
    \end{array}
  \]
\end{ClozeArea}
\noindent\textbf{Rückwärtseinsetzen:}
\begin{itemize}[leftmargin=1.4em,itemsep=0.55em,topsep=0.3em]
  \item in Zeile 2 einsetzen:
    \begin{ClozeArea}[32]{3}
      \[
        (x_B-x_A)m=y_B-y_A
        \quad\Longleftrightarrow\quad
        m=\frac{y_B-y_A}{x_B-x_A}
        \qquad (x_A\ne x_B).
      \]
    \end{ClozeArea}
  \item in Zeile 1 einsetzen:
    \begin{ClozeArea}[32]{3}
      \[
        b + m x_A = y_A \quad \Longleftrightarrow \quad b=y_A-m x_A
      \]
    \end{ClozeArea}
\end{itemize}
\begin{KeyPointBox}{Geradengleichung durch zwei Punkte}{real_functions_two_point_form}
    Die Gerade, die durch die Punkte $P_A=(x_A,y_A)$ und
    $P_B=(x_B,y_B)$ mit $x_A\ne x_B$ verläuft, hat die Funktionsgleichung
    \[
      g(x)=\textcolor{THMBlue}{m}x+\textcolor{green!55!black}{b}
    \]
  \begin{minipage}[c]{0.52\linewidth}
    mit ...
    \begin{itemize}[leftmargin=1.3em,itemsep=0.4em,topsep=0.3em]
      \item der Steigung:
        \[\textcolor{THMBlue}{m=\frac{y_B-y_A}{x_B-x_A}} \qquad \text{\textcolor{THMBlue}{(vgl. Steigungsdreieck)}}.\]
      \item dem $y$-Achsenabschnitt:
        \[\textcolor{green!55!black}{b=y_A-mx_A \qquad \text{oder} \qquad b=y_B-mx_B}.\]
    \end{itemize}
  \end{minipage}\hfill
  \begin{minipage}[c]{0.46\linewidth}
    \centering
    \CachedTikzFigure{content/resources/chapters/real_functions/tikz/functions_two_point_line.tex}
  \end{minipage}
\end{KeyPointBox}
\begin{ExampleBox}{Gerade durch zwei Punkte}{real_functions_two_point_numeric_example}
  Welche Gerade verläuft durch $P_A=(1,1)$ und $P_B=(3,2)$?
  \par\smallskip
  \begin{minipage}[c]{0.48\linewidth}
    Aus den beiden Punkten bestimmen wir:
    \begin{itemize}[leftmargin=1.3em,itemsep=0.35em,parsep=0pt,
      topsep=0.3em,after=\vspace{0.65em}]
      \item die \textcolor{THMBlue}{Steigung}:
        \begin{ClozeArea}[15]{3}
          \[
            \textcolor{THMBlue}{m=\frac{y_B-y_A}{x_B-x_A}
              =\frac{2-1}{3-1}=\frac12},
          \]
        \end{ClozeArea}
      \item den \textcolor{green!55!black}{$y$-Achsenabschnitt}:
        \begin{ClozeArea}[15]{3}
          \[
            \textcolor{green!55!black}{b=y_A-mx_A
              =1-\frac12\cdot1=\frac12}.
          \]
        \end{ClozeArea}
    \end{itemize}
    Damit lautet die Geradengleichung
    \begin{ClozeArea}[16]{3}
      \[
        g(x)=\textcolor{THMBlue}{\frac12}x+\textcolor{green!55!black}{\frac12}.
      \]
    \end{ClozeArea}
  \end{minipage}\hfill
  \begin{minipage}[c]{0.51\linewidth}
    \raggedleft
    \CachedTikzFigure{content/resources/chapters/real_functions/tikz/functions_two_point_numeric_example.tex}
  \end{minipage}
\end{ExampleBox}
\begin{ExampleBox}{Kalibrierung eines Wegsensors}{real_functions_sensor_calibration}
  Ein Wegsensor dient dazu, eine Strecke mithilfe einer Spannung zu messen.
  Dazu benötigen wir den Zusammenhang $U(s)$ zwischen Weg $s$ und Spannung
  $U$, den wir als affine Abbildung $U(s)=ms+b$ annehmen.

  \par\smallskip
  Zur Kalibrierung messen wir die Spannung bei zwei bekannten Wegen. Bei
  $s_A=10\,\mathrm{mm}$ erhalten wir $U_A=1\,\mathrm{V}$ und bei
  $s_B=50\,\mathrm{mm}$ erhalten wir $U_B=9\,\mathrm{V}$.\\
  \par\smallskip
  \begin{minipage}[c]{0.48\linewidth}
    Aus den beiden Messpunkten bestimmen wir:
    \begin{itemize}[leftmargin=1.3em,itemsep=0.35em,parsep=0pt,
      topsep=0.3em,after=\vspace{0.65em}]
      \item die \textcolor{THMBlue}{Steigung}:
        \begin{ClozeArea}[15]{7}
          \[
            \textcolor{THMBlue}{\begin{aligned}
              m&=\frac{\Delta U}{\Delta s}=\frac{U_B-U_A}{s_B-s_A}\\[0.65em]
               &=\frac{9\,\mathrm{V}-1\,\mathrm{V}}
                      {50\,\mathrm{mm}-10\,\mathrm{mm}}
                 =0{,}2\,\frac{\mathrm{V}}{\mathrm{mm}}.
            \end{aligned}}
          \]
        \end{ClozeArea}
      \item den \textcolor{green!55!black}{$y$-Achsenabschnitt}:
        \begin{ClozeArea}[15]{7}
          \[
            \textcolor{green!55!black}{\begin{aligned}
              b&=U_A-\textcolor{THMBlue}{m} \, s_A\\[0.65em]
               &=1\,\mathrm{V}
                 -\textcolor{THMBlue}{0{,}2\,\frac{\mathrm{V}}{\mathrm{mm}}}
                  \cdot10\,\mathrm{mm}\\[0.65em]
               &=-1\,\mathrm{V}.
            \end{aligned}}
          \]
        \end{ClozeArea}
    \end{itemize}
  \end{minipage}\hfill
  \begin{minipage}[c]{0.51\linewidth}
    \raggedleft
    \CachedTikzFigure{content/resources/chapters/real_functions/tikz/functions_sensor_calibration.tex}
  \end{minipage}
  Damit lautet die gesuchte Kennlinie
  \begin{ClozeArea}[17]{3}
    \[
      U(s)=\textcolor{THMBlue}{0{,}2\,\frac{\mathrm{V}}{\mathrm{mm}}} \cdot s \textcolor{green!55!black}{-1\,\mathrm{V}}.
    \]
  \end{ClozeArea}
\end{ExampleBox}
